% Section 10.3, from lectures/L10/README.md: the objectives, the evaluation questions, and the course
% review.

\section{Review}
\label{c10:sec:review}

\needspace{6\baselineskip}
\subsection*{What you should be able to do now}
After this chapter, you should be able to:
\begin{itemize}
  \item \textbf{Write \code{app::EchoNode}} against the driver interface and host-test it on the
    UART stub of \chapref{6} with no hardware, then bring it up as the final rung of the ladder.
  \item \textbf{Wire the two-link bench correctly}, level shifter included, and explain what each
    SPI line and each UART line carries.
  \item \textbf{Climb the bring-up ladder methodically} and state what each rung establishes that
    the previous one cannot.
  \item \textbf{Explain what an integration test catches that no host test can}, and where the top
    of the test pyramid sits for this system.
\end{itemize}

\needspace{6\baselineskip}
\subsection*{Questions to test yourself}
You should be able to explain:
\begin{enumerate}
  \item The real-data-plane rung (d) works but the EchoNode rung (e) does not: which layer is
    implicated, and which layers has the ladder already cleared?
  \item Rung c loops \code{tx} back to \code{rx} on the board and passes, yet rung d fails against
    the terminal: what does that isolate, given that the peripheral's own datapath is common to
    both?
  \item What does the control-plane rung (b) prove that a passing \code{uart_top_tb} in simulation
    did not?
  \item If one SPI line is left unshifted, wired 5\,V straight to a 3.3\,V pin, what is the likely
    symptom, and which rung exposes it first?
\end{enumerate}

\needspace{6\baselineskip}
\subsection*{Course review}
The pattern you have just completed --- a peripheral in VHDL, a driver in C++, one shared register
contract, integrated across a real chip boundary --- is the whole FPGA-meets-MCU method. The same
method scales to far harder protocols, and the SPI transport this book handed you as a black box is
itself built from the wire up two courses later, in \emph{SPI: The MCU-FPGA Transport, from the Wire
Up}.
